\begin{lemma}[Step 1, Subcase 2.2 at a PE zero divider]
Let $G,e_0,[k],T,\varphi$ be a finite loopless multigraph with exact edge
identities, distinguished uncolored edge, color set, positive-rung
certified ETT, and generating coloring.  Assume maximum property,
properness, exact one-uncoloredness, and the complete six-part lower-rung
package.  Let $H$ be a good hierarchy over a hierarchy base, with
$d<q_0$ dividers.  Assume the same-rung fewer-than-$q_0$ elementarity
hypothesis for arbitrary certified inputs and generating colorings, the
large bound $|\overline\varphi(T_n)|\geq2n+11$, the separate terminal
certificate
\[
 T\text{ is }
 \left(\bigcup_{\eta\in D^H_{n,d+1}}\Gamma^d_\eta\right)^-
 \text{-closed under }\varphi,
\]
and elementarity of $V(T)$.  Put $A=H.\mathit{starred}(d)$.

Let $\alpha\ne\beta$, with $\alpha$ missing and singleton-closed on $T$
and $\beta$ not missing on $T$.  Let $P_1,P_2,P_3$ be pairwise
support-disjoint exact $(T,\varphi,\{\alpha,\beta\})$-exit paths.  Their
outside ends $a_1,a_2,a_3$ and tree ends $b_1,b_2,b_3$ are pairwise
distinct, with $b_1\prec_T b_2\prec_T b_3$.  Retain the exact
$\beta$-colored boundary-edge identity $f_j$ of each path incident with
$b_j$.  Suppose $b_3\in V(A)$.

Assume that a selected color $\theta\ne\alpha,\beta$ is missing and
singleton-closed on $A$ and singleton-closed on $T$.  Assume also the
explicit no-smaller-presentation consequence of the global minimum: for
every proper, exactly one-uncolored coloring $\sigma$ stable from
$\varphi$ on the final ladder tree, arbitrary replacement input,
corresponding same-rung maximum ETT $T'$, good hierarchy with fewer than
$d$ dividers, its own large estimate, and its separate terminal
minus-closed certificate, there do not exist two support-disjoint exact
exit legs with distinct outside ends for a missing singleton-closed first
color.  Assume additionally that $d=0$ and that the terminal extension
kind is PE.  Then the displayed Subcase 2.2 data are contradictory.
\end{lemma}
\begin{proof}
Write $n$ for the rung number, $T_n$ for the last ladder tree, $D_i$ for
the exact cumulative stored-color set at slot $i$, and $\varphi_i$ for
the exact coloring at that slot.  Thus $\varphi=\varphi_n$ by the
generating-coloring field of \noderef{ExtendedTashkinovTree}.  We first
make the three exit legs used in every branch completely explicit.

The starred tree $A$ is an ordered initial segment of $T$.  Unfold
\noderef{TashkinovHierarchy}.  If $d>0$, then
\noderef{HierarchyStarredStage} identifies $A$ with the $d$-th hierarchy
stage, and repeated use of the nestedness field of the hierarchy gives
the required initial-segment map to its terminal stage $T$.  If $d=0$,
the same conclusion follows from the core-inside-first field and then the
nestedness and terminal-stage fields of the hierarchy.  All these are
entrywise statements by \noderef{TreeSequenceInitialSegment}.
Unfolding \noderef{TreeVertexOrder} in the two displayed order
relations gives the occurrences of $b_1,b_2,b_3$ at strictly increasing
positions of $T$.  Since $b_3\in V(A)$ and the entries of $A$ form an
initial block of those of $T$, the occurrence of $b_3$ is in that block.
The two earlier occurrences are consequently in the same block, so
$b_1,b_2,b_3\in V(A)$ by \noderef{TreeSequenceVertexSet}.

Unfold each given \noderef{TreeExitPath}.  The edge-node immediately
preceding the terminal vertex-node $b_j$ represents an exact boundary
edge of $T$: its other adjacent vertex-node on the path is outside $T$,
because $b_j$ is the only vertex-node of the path in $V(T)$.  By
\noderef{TwoColorIncidenceGraph} that edge has color $\alpha$ or
$\beta$.  It cannot have color $\alpha$, since
\noderef{TreeColorSetClosed} says that the exact $\alpha$-colored
boundary set is empty.  Hence this last boundary edge is
$\beta$-colored.  The separately retained identity $f_j$ has the same
two properties by hypothesis; unfolding \noderef{BoundaryColorEdges}
and \noderef{BoundaryEdges} shows in particular that its end other than
$b_j$ lies outside $T$.  Since $b_j\in V(A)\subseteq V(T)$, the same
identity $f_j$ is also a $\beta$-colored boundary edge of $A$.

We next justify which $\alpha,\theta$ components may be exchanged.  If an
$\alpha$-edge node occurs on $P_j$, it is not incident with $T$:
otherwise one of its incident vertex-nodes on this alternating path
would be in $T$, hence would be $b_j$, and the edge would be an
$\alpha$-colored boundary edge, contrary to the preceding closedness.
More generally, its full \noderef{TwoColorComponent} for
$\alpha,\theta$ contains no vertex of $T$.  For if it did, take a
shortest incidence-graph walk from the external edge to a vertex-node in
$T$ and its first vertex-node in $T$.  The preceding exact edge identity
crosses the boundary of $T$ and, by
\noderef{TwoColorIncidenceGraph}, has color $\alpha$ or $\theta$.
This contradicts respectively the assumed singleton closedness of
$\alpha$ or of $\theta$ on $T$.  The identical first-crossing argument
shows that every $\alpha,\theta$ component containing an outside vertex
is wholly outside $T$.

Exchange $\alpha$ and $\theta$ on every full $\alpha,\theta$ component
outside $T$, and call the result $\nu$.  Finiteness makes this a finite
choice of roots, and the preceding paragraph verifies the outside-support
clause of \noderef{OutsideKempeRecoloring}.  By
\noderef{OutsideKempeRecoloringStable}, $\nu$ is proper, leaves exactly
$e_0$ uncolored, fixes every exact edge incident with $T$, preserves the
missing-color set at each vertex of $T$, and is
$(T_n,C,\varphi_n)$-stable for every chosen control set $C$: after
unfolding \noderef{StableColoring}, restrict its incident-edge and vertex
clauses from $T$ to the ordered initial segment $T_n$.

For $j=1,2,3$, form a walk $Q_j$ with the same ordered vertex- and
edge-nodes as $P_j$.  Each old $\beta$-edge remains $\beta$ because
$\beta\ne\alpha,\theta$.  Every old $\alpha$-edge lies in one of the
selected outside components just proved and becomes $\theta$.  Thus
every successive pair of nodes of $P_j$ is adjacent in
$I_G(\nu;\theta,\beta)$, so $Q_j$ is a walk there; it is the same simple
node list and has the same positive length and support as $P_j$.
Its outside end still misses one of its two colors.  Explicitly, if
$\beta$ was missing at $a_j$, then the definition
\noderef{MissingColors}, together with \noderef{PresentColors}, says
that no incident edge had color $\beta$; none acquires that color in an
$\alpha,\theta$ exchange, so $\beta$ remains missing.  If instead
$\alpha$ was missing, then no incident edge had color $\alpha$.  The
whole $\alpha,\theta$ component through the outside vertex $a_j$ was
selected, so the $\theta$-edges incident with $a_j$ become $\alpha$ and
the new $\theta$-edges are exactly its old $\alpha$-edges; hence
$\theta$ is missing there under $\nu$.  We already proved that all three
tree ends lie in $A$, while $a_j\notin V(T)$ implies
$a_j\notin V(A)$.  Finally, the only vertex-node of $Q_j$ in $A$ is
$b_j$, because the corresponding claim for $P_j$ excludes every other
vertex even from the larger tree $T$.  This checks every field of
\[
       \operatorname{TreeExitPath}(A,\nu,\theta,\beta,a_j,b_j,Q_j).
       \tag{1}
\]
The supports remain pairwise disjoint, the three outside ends remain
pairwise distinct, and no claim about separation of the three
transformed full components has been made.

Apply \noderef{HierarchyLEMStable} to the just constructed outside
recoloring.  It supplies a corresponding ETT $\mathcal T^\nu$ displayed
on $T$, a hierarchy base $B^\nu$, a hierarchy $H^\nu$, and a good
hierarchy on it.  The corresponding certificate retains the rung, all
ordered ladder trees, exact connecting-edge sets, stored sets, and kinds;
the same-hierarchy conclusion retains every ordered stage, divider,
stored-color set, and Gamma pair.  The result also supplies maximum
property and transports the given terminal minus-closed certificate.
The stability established above identifies the missing-color sets on
$T_n$ under $\nu$ and $\varphi_n$, so the large bound is unchanged.

By the zero-divider hypothesis, $d=0$.  By \noderef{HierarchyStarredStage}, $A$ is the
actual hierarchy core.  We treat the two possible terminal constructions
separately.

The terminal extension is PE by hypothesis.  Use the transported
corresponding ETT $\mathcal T^\nu$ and its transported base $B^\nu$, not
the obsolete pre-recoloring PE data.  Its PE field gives an actual
support vertex $x$ in $T_n$, an exact closure $R_n$ of the prefix through
$x$, and identifies its core with the ordered union $T_n\vee R_n=A$.
The prefix contains $x$ by the entrywise definition
\noderef{TreeSequencePrefix}.  The augmentation sequence in
\noderef{TashkinovClosure}, after unfolding
\noderef{TashkinovAugmentationSequence} and
\noderef{TashkinovAugmentation}, consists of ordered initial-segment
steps.  Induction through those steps therefore retains the support
prefix in $R_n$, so $x\in V(R_n)$ as well.  Choose a color $c$ missing
at $x$ by the same $k\geq\Delta(G)+1$ degree count as above.  Thus
\noderef{TreeMissingColors} puts $c$ in both
$\overline\nu(T_n)$ and $\overline\nu(R_n)$.

The pre-step elementarity required next is available for these
transported data.  If $n>1$, apply \noderef{CertifiedETTPredecessor} to
$\mathcal T^\nu$ and then lower assertion (i); if $n=1$, apply the
elementarity conclusion of \noderef{TechnicalTheoremBase} to the exact
initial tuple of the replacement series.  Therefore
\noderef{ParallelCommonMissingClosed}, with $\mathcal T^\nu$ and
$B^\nu$, makes every common missing color, and in particular $c$,
singleton-closed on the exact core $A$.

The retained identity $f_3$ is still a $\beta$-colored boundary edge of
$A$: its inside endpoint is $b_3\in V(A)$, its other endpoint is outside
$T$ and hence outside $A$, and the first outside recoloring fixed its
color.  Consequently $c\ne\beta$.  Moreover $\beta$ is present at $x$.
For if it were missing there, then $x\in V(T_n)\cap V(R_n)$ would put
$\beta$ in both tree-missing sets; a second application of
\noderef{ParallelCommonMissingClosed} would make $A$
singleton-$\beta$-closed, contradicting the exact boundary identity
$f_3$.  Unfolding \noderef{MissingColors} and
\noderef{PresentColors} therefore supplies an exact $\beta$-colored edge
incident with $x$.

If $\theta=c$, put $\rho=\nu$; otherwise exchange $\theta,c$ on all
full components outside $A$.  The first outside recoloring fixes incident
edges and missing sets on $A$, so the original $\theta$-missingness and
singleton closedness on $A$ hold under $\nu$; singleton-$c$-closedness
was just proved.  The first-crossing argument again verifies the outside
support required by \noderef{OutsideKempeRecoloring}.  In the
nonidentity case, apply \noderef{OutsideKempeRecoloringStable} first to
$T_n$ with control set $D_n$ and then to $R_n$ with empty control;
components outside $A$ are outside each constituent.  In the identity
case use reflexivity of stability.  The two resulting clauses are
exactly \noderef{CombinedStableColoring} relative to the transported
generating coloring $\nu$.  They also give properness, exact
one-uncoloredness, and preserve the missing sets at every vertex of $A$.

In particular $c$ remains missing at $x$, hence belongs to
$\overline\rho(A)$.  Every exact edge incident with $A$ is fixed, so the
empty $c$-colored boundary under $\nu$ remains empty under $\rho$; thus
$A$ is singleton-$c$-closed under $\rho$.  The exact $\beta$-edge at
$x$ is fixed as well, so $\beta$ remains present there.  Repeating the
edge-by-edge and outside-end argument used to construct (1) gives three
exact $(A,\rho,\{c,\beta\})$-exit paths with the old pairwise-disjoint
supports and pairwise-distinct outside ends.

The leaf argument above now applies verbatim in
$I_G(\rho;c,\beta)$.  It certifies, via
\noderef{TwoColorPathComponent}, the full component through $x$ and the
full components through the three outside ends as path components; the
first meets $A$ at $x$, and each latter component meets $A$ along its
exit path.  The spanning-path argument chooses one outside component
distinct from the component through $x$, again without asserting that
the three transformed leg components are mutually distinct.  Apply
conclusion (iii) of \noderef{interchange} to $\mathcal T^\nu$, its PE
base $B^\nu$, the just constructed
\noderef{CombinedStableColoring}, and the proved $c$-missing and
singleton-$c$-closed properties on $A$.  It says that $c,\beta$ are
$A$-interchangeable.  The definition \noderef{TreeInterchangeable} then
equates the two distinct full path components meeting $A$, the final
contradiction.  This completes Subcase 2.2.
\end{proof}
